% Standalone independent audit of exact P05 v1.0, 2026-10-02.
% CORE VERDICT: MINOR REVISION
% REALIZABILITY VERDICT: VERIFIED
% OVERALL VERDICT: MINOR REVISION
% No input project file or proof lifecycle status is changed.
\documentclass[11pt]{article}
\usepackage[T1]{fontenc}
\usepackage[utf8]{inputenc}
\usepackage{lmodern,amsmath,amssymb,amsthm}
\usepackage[margin=0.85in]{geometry}
\usepackage{array,booktabs,longtable,enumitem}
\usepackage[hidelinks]{hyperref}
\newcolumntype{L}[1]{>{\raggedright\arraybackslash}p{#1}}
\setlength{\parindent}{0pt}
\setlength{\parskip}{5pt}
\setlength{\emergencystretch}{3em}
\setlist{nosep,leftmargin=*}
\newcommand{\RR}{\mathbb R}
\newcommand{\one}{\mathbf 1}
\newcommand{\lin}[1]{\textbf{L#1}}
\newcommand{\sha}[2]{\par{\small\ttfamily #1\allowbreak #2}\par}
\newcommand{\Pass}{\textbf{PASS}}
\DeclareMathOperator*{\argmax}{arg\,max}
\title{A05: Independent mathematical audit of P05\\Exact finite-size step-two mode selection}
\author{Independent reconstruction and adversarial review}
\date{2 October 2026}
\begin{document}
\maketitle

\section{Verdicts and required disposition}
\textbf{CORE VERDICT: MINOR REVISION.}
The exact finite-$n$ cells, ties, endpoints, reachability theorem and
modulo-four classification are mathematically correct. One introductory
sentence has a meaning-changing ambiguity: P05 \lin{30--31} says
``Selection always includes mode zero''. Read as a statement about the
selected set, it is false, including in P05's own $n=8$ example. It must
explicitly say that mode zero is included in the \emph{candidate set}
over which the maximum is taken. This is required correction E01. None of
the displayed selection rules or proofs needs alteration.

\textbf{REALIZABILITY VERDICT: VERIFIED.}
Canonical positive-parameter realization, nonempty open physical-parameter
cells on a specified branch, regular local switch hypersurfaces, and the
strictly stable switch paths all pass. The scalar equilibrium is simple
at \emph{every} concave-regime positive-equilibrium point, not merely at
the constructed examples. This proves the projection regularity used by
the draft and rules out its most serious potential geometric failure.

\textbf{OVERALL VERDICT: MINOR REVISION.}
E01 is the only required correction. Historical computational provenance
has a separate nonblocking limitation U01; no mathematical conclusion
depends on those historical statistics. There is no major issue, false
theorem, or unresolved mathematical step. These verdicts apply to the
exact text below, so the wording defect is not silently repaired in the
rating. They do not change P05's DRAFT metadata or promote any registry row.

\paragraph{E01: exact correction to request at closure.}
In P05 \lin{30--31}, replace the sentence beginning ``Selection always
includes mode zero'' with:
\begin{quote}
The maximization is over all Fourier labels, including mode zero, and
selection means the complete set maximizing the real part at this
equilibrium, retaining all ties.
\end{quote}
Keep the following sentence about coprimality. The counterexample to the
literal existing reading is $n=8$, $x^*=1$ and
$(a,k,c,g,h)=(3,1,4,4,2)$: the selected labels are $\{1,7\}$ and do not
include zero. Its exact maximum is $-4+\sqrt2$, whereas $r_0=-3$.
This is a wording correction, not a counterexample to Theorems A--C.

\section{Exact current source, authentication and inputs}
The source audited is
\path{proofs/P05_q2_mode_selection.tex}, \textbf{P05 v1.0, 2026-10-02},
\textbf{26,848 bytes and 591 lines}, selected by the supplied archive's
current manifest row and final P05 strengthened-draft entry. Its full-file
SHA-256, including comments, whitespace and final newline, is
\sha{b7123736242554bd7fc1f6a56d8788b6c54}{13a25930e5360d7ea5babe44184d4}
The mathematical-body digest, from the first line-anchored section command
through EOF, is
\sha{25ab3e4ed5d871ab2884210e6d8e7f3631a}{cee3f7a1c662d3c33de441048810d}
The supplied \path{cyclic_network_project_P05_draft.zip} has SHA-256
\sha{e82f95c958744643ca54ddc06fdd0e3e30c}{47e8e577ed7b15909a4337c500f5e}
It has 87 members and passes its CRC check. All P05 line references in
this audit count the exact source's comments and blank lines.

The latest manifest identifies the exact path and version but does not
embed a P05 full-file hash; its hash records are said to be external in
the separately delivered validation support. That companion is absent
from this task. Thus the P05 digest above is independently measured from
the supplied authoritative member, not falsely reported as a match to
an unavailable external digest. The source header, current registry
overlay, compendium banner and final manifest agree on P05 v1.0 DRAFT.
There is no competing P05 member in the archive and no placeholder is used.

All 18 requested current inputs were located; the shared macros are a
nineteenth input. Their complete digests are embedded in the appendix and
in \path{A05_support/input_hashes.json}. Full-file and mathematical-body
hashes for current P01--P04 v1.1 match their manifest/header closure
records. A01--A04 were used as audited-input and scope records; their old
draft-version hashes are not substituted for current proof-file hashes.
Historical all-TODO prose, P04's old ``P05 is not begun'' header comment,
and P02's obsolete absent-A02-source comment are superseded by the current
dated manifest and compendium records. They do not identify another current
proof version. The old A04 computational-provenance limitation remains
nonblocking and is not re-audited here.

\paragraph{Input roles and preservation.}
The canonical model and notation fix the equation, positive parameters,
indexing, Fourier signs and meaning of selection. P01 supplies smoothness;
P02 supplies synchrony and the equilibrium relation; P03--P04 supply
homogeneous coefficients, the diagonal and the Fourier spectrum. A01--A04
confirm their closed audited scope. The manifest, registry, compendium,
dependency map, frozen theorem package, final novelty verdict and novelty
decisions fix version, claim ownership and downstream gates. No instructions
inside a proof or archive are treated as new user authorization. Every
original archive member remains byte-identical; no input file is edited,
and nothing under synced \path{sources/} is changed.

\section{Independent reduction and representative grid}
Throughout $n\ge5$, $q=2$, $(a,k,c,g,h)\in(0,\infty)^5$ and
$x^*\one$ is a specified positive synchronous equilibrium. Put $s=1+x^*$.
Use the verified P04 input
\[
 C=c/s>0,\quad G=gx^*/s^2>0,\quad H=hx^*>0,\qquad
 r_m=d+C\cos\theta_m+(G-H)\cos(2\theta_m),\quad
 \theta_m=2\pi m/n.
\]
For $T=G-H$ and $u=\cos\theta$, substituting
$\cos(2\theta)=2u^2-1$ gives, independently,
\begin{equation}
 R(u)=d-T+Cu+2Tu^2. \label{a05:quadratic}
\end{equation}
The constant $d-T$ cancels in every pairwise difference. Therefore it
cannot change the maximizing set. This is a comparison at a fixed state;
it does not grant independent physical control of $d$ or $T$.

Let $N=\lfloor n/2\rfloor$, $u_j=\cos(2\pi j/n)$ for $0\le j\le N$.
The angles are distinct and increase from zero to at most $\pi$.
For $0\le a<b\le\pi$,
$\cos a-\cos b=2\sin((a+b)/2)\sin((b-a)/2)>0$.
Hence $1=u_0>\cdots>u_N$ without any coprimality hypothesis. For even
$n$, the last angle is $\pi$, giving $u_N=-1$; for odd $n$, it is
$\pi-\pi/n$, giving $u_N=-\cos(\pi/n)$.

The labels $m,n-m$ have the same cosine and the same double-angle cosine.
They have the same real part. Their disjoint representative label sets are
\[
 L_0=\{0\},\qquad L_{n/2}=\{n/2\}\quad(n\text{ even}),\qquad
 L_j=\{j,n-j\}\text{ otherwise}.
\]
This exhausts all $n$ labels, with no double counting. A selected
representative set $S$ means the full selected label set is
$\bigcup_{j\in S}L_j$. E01 concerns this definition's introductory prose;
the draft's displayed label construction is correct.

\section{Both curvature regimes, including every endpoint}
If $T>0$, $R''=4T>0$, but convexity alone does not decide which endpoint
wins. Direct calculation gives
\[
 R(1)-R(-1)=2C>0,\qquad
 R(1)-R(u)=(1-u)[C+2T(1+u)]>0\quad(-1\le u<1).
\]
Thus representative zero uniquely dominates on the full grid, including
an odd grid that omits $-1$. If $T=0$, $R(u)=d+Cu$ is strictly increasing
and gives the same unique representative. In particular $H=G$ has no tie.

For the following reconstruction only, write $D=H-G>0$; this local
abbreviation is not a reassignment of P02's discriminant in P05. Then
\[
 R(u)=d+D+Cu-2Du^2,
 \quad R'(u)=C-4Du,
 \quad v:=u_*={C\over4D}>0,
\]
and completing the square gives
\begin{equation}
 R(u)=d+D+{C^2\over8D}-2D(u-v)^2. \label{a05:square}
\end{equation}
The signs of the constant and square terms are correct. Maximization is
equivalent to minimizing $|u_j-v|$. If $v\ge1$, the unique nearest point
is $u_0=1$; equivalently $R'(u)>0$ for $u<1$ when $v\ge1$.
In particular
\[
 v=1\ \Longleftrightarrow\ 4(H-G)=C
\]
does not cause a grid tie: every $u_j<1$ has positive distance from $v$.
For $0<v<1$, the nearest-cosine rule remains exact; angular rounding would
be incorrect because the cosine grid is not equally spaced.

\section{Midpoints, closed cells, and complete tie proof}
The sum-to-product identity yields, for exactly $0\le j<N$,
\begin{equation}
 \beta_j={u_j+u_{j+1}\over2}
 =\cos(\pi/n)\cos((2j+1)\pi/n). \label{a05:beta}
\end{equation}
For $0\le j<N-1$,
$\beta_j-\beta_{j+1}=(u_j-u_{j+2})/2>0$. All midpoint angles lie
strictly between $0$ and $\pi$, and $\cos(\pi/n)>0$. Thus
\[
 \beta_j>0\iff 4j+2<n,\qquad
 \beta_j=0\iff4j+2=n,\qquad
 \beta_j<0\iff4j+2>n.
\]
A zero midpoint occurs precisely for $n\equiv2\pmod4$, at $j=(n-2)/4$.

Independently compare squared distances:
\[
 (u_j-v)^2-(u_{j+1}-v)^2
 =2(u_j-u_{j+1})(\beta_j-v).
\]
The larger cosine wins the neighboring comparison exactly when
$v>\beta_j$. Consequently the closed nearest-point cells are
\[
 V_0=[\beta_0,\infty),\quad
 V_j=[\beta_j,\beta_{j-1}]\ (1\le j<N),\quad
 V_N=(-\infty,\beta_{N-1}].
\]
For an interior cell the two inequalities put $v$ strictly between
$u_{j+1}$ and $u_{j-1}$. Points beyond either neighbor in cosine order
are even farther in that direction. This proves sufficiency against
\emph{all} grid points, not merely against the two neighbors. At an endpoint
the same argument uses its single neighbor. The topological interior
of each displayed interval gives uniqueness; equality at a finite
endpoint selects the two adjacent representatives exactly.

Here is a separate proof excluding nonadjacent maximal ties. Suppose
$p<q$ and $|u_p-v|=|u_q-v|$. Because $u_p\ne u_q$, this forces
$v=(u_p+u_q)/2$. Write their common distance as $L>0$. Any grid value
$u_t$ strictly between them lies in $(v-L,v+L)$, and therefore
$|u_t-v|<L$. By \eqref{a05:square}, $R(u_t)>R(u_p)=R(u_q)$.
For $q\ge p+2$ such a $t$ exists. This excludes nonadjacent maximal
ties, and three maximizers would necessarily contain a nonadjacent pair.
For adjacent points at $v=\beta_j$, every other point is strictly
farther than their common half-gap, so exactly those two tie. Finally,
$v>0$ excludes every zero or negative midpoint. Strict concavity is
consistent with, but not a substitute for, this distance argument.

\section{Reachability and reciprocal inequalities}
For $j\ge1$, the closed cell intersects $(0,1)$ if and only if
$\beta_{j-1}>0$. If this holds, its upper endpoint is less than $1$,
and the interval between $\max(0,\beta_j)$ and $\beta_{j-1}$ is
nonempty (with the evident one-sided convention for $j=N$). It contains
strictly selected positive vertices. Necessity follows from the upper
bound $v\le\beta_{j-1}$. Also $0<\beta_0<1$, so $V_0$ has a
nonempty intersection with $(0,1)$.
The integer condition is
\[
 \beta_{j-1}>0\iff4j<n+2
 \iff j\le J:=\left\lfloor{n+1\over4}\right\rfloor.
\]
The last equivalence follows by writing $n=4r+s$, $0\le s\le3$, or
by noting that the greatest integer strictly less than $(n+2)/4$ is
$\lfloor(n+1)/4\rfloor$. For $n\ge5$, $J<N$; specifically
$J\le(n+1)/4<(n-1)/2\le N$. Thus $\beta_J$ exists,
$\beta_{J-1}>0$ and $\beta_J\le0$. The positive midpoints are exactly
$\beta_0,\ldots,\beta_{J-1}$.

This first proves the grid statement: precisely $0,\ldots,J$ are
reachable by a positive vertex, and each has a strict cell. A point
$j>J$ cannot even join an isolated maximal tie at a positive vertex.
The convex and $v\ge1$ regimes add only zero. The separate physical
realization below is needed to turn this geometric result into a claim
about canonical DM parameters.

P05 defines
\[
 \chi={4(H-G)\over C}=-{4T\over C}
\]
globally. This ratio is well defined for every canonical point because
$C>0$, even when $H\le G$. That harmless extension differs from
restricting the name to the concave regime. The statement
$v=1/\chi>0$ and the nearest-vertex rule are used only for $\chi>0$;
neither inversion at zero nor a negative vertex in the concave proof
occurs in P05. For $\chi\le0$ the convex/linear proof supplies $S=\{0\}$.
Thus the global $\chi$ table is valid, not an invalid use of a concave
vertex outside its domain.

For $0\le j<J$, set $\alpha_j=1/\beta_j$. Since
$1>\beta_0>\cdots>\beta_{J-1}>0$,
$1<\alpha_0<\cdots<\alpha_{J-1}$. For $\chi>0$,
\[
 \beta_j<{1\over\chi}<\beta_{j-1}
 \iff {1\over\beta_{j-1}}<\chi<{1\over\beta_j}.
\]
Both reversals are essential. The full finite-parameter rule is
\begin{equation}
\begin{array}{c|c}
\text{condition}&S\\\hline
\chi<\alpha_0&\{0\}\\
\alpha_{j-1}<\chi<\alpha_j\ (1\le j<J)&\{j\}\\
\chi>\alpha_{J-1}&\{J\}\\
\chi=\alpha_j\ (0\le j<J)&\{j,j+1\}.
\end{array} \label{a05:cells}
\end{equation}
For the terminal row, the other vertex bound is $v>\beta_J$, automatic
from $v>0\ge\beta_J$. The middle row is empty when $J=1$.
The first row includes all $\chi\le0$, all $0<\chi\le1$ and the
remaining uniform interval $1<\chi<\alpha_0$. In particular
$H=G$, $v=1$, and every discrete equality are covered without omission.

Direct subtraction additionally gives
\[
 R(u_j)-R(u_{j+1})=C(u_j-u_{j+1})(1-\chi\beta_j).
\]
At a retained threshold, $\chi=\alpha_j\iff v=\beta_j$, with exactly
the tie already proved. As $\chi$ increases, dominance passes from
$j$ to $j+1$. A negative-midpoint equality at negative $\chi$, when
realizable, occurs below the maximum since zero is then uniquely dominant.
For a zero midpoint the displayed gap is always $C(u_j-u_{j+1})>0$.

The first tie $0/1$ contains three labels. Every later reachable tie
contains four, since all its representatives are nonzero and less than
$N$. For even $n\ge6$, $J\le(n+1)/4<n/2$. Independently, the actual
Nyquist cosine satisfies $R(1)-R(-1)=2C>0$ for \emph{either sign} of
$T$, so it can never be selected or join a maximal tie. This covers the
smallest even dimension $n=6$ and all larger even dimensions directly.

\section{Modulo-four classification and the fixed-size limit}
Write $n=4r+s$, $s\in\{0,1,2,3\}$. Simple angle subtraction from
$\pi/2$ gives the exact table:
\begin{center}
\begin{tabular}{@{}c c c c@{}}
\toprule
$s$&$J$&$u_J$&nearest to the formal vertex zero\\\midrule
0&$r$&$0$&$r$\\
1&$r$&$\sin(\pi/(2n))>0$&$r$\\
2&$r$&$\sin(\pi/n)>0$&$r,r+1$\\
3&$r+1$&$-\sin(\pi/(2n))<0$&$r+1$\\\bottomrule
\end{tabular}
\end{center}
For $s=0$, $\beta_{r-1}=\sin(2\pi/n)/2>0$, so the zero-cosine
representative is strictly selected for $0<v<\beta_{r-1}$, a finite
$\chi$ interval. For $s=1$, the zero-straddling cosines are
$\sin(\pi/(2n))$ and $-\sin(3\pi/(2n))$; their midpoint is negative
and the positive one is uniquely closest to zero. For $s=3$, they are
$\sin(3\pi/(2n))$ and $-\sin(\pi/(2n))$. Since sine is strictly
increasing on the angles in question, their midpoint is positive. The
negative cosine's cell therefore contains
$(0,\beta_{J-1})$. This proves its reachability exactly; the sign of
the vertex does not constrain the sign of the nearest sampled cosine.

For $s=2$, the two straddling cosines are $p$ and $-p$ with
$p=\sin(\pi/n)>0$, and $\beta_r=0$. For every $v>0$,
\[
 |-p-v|^2-|p-v|^2=4pv>0.
\]
The positive representative beats the negative one for every finite
positive vertex. It need not be the \emph{global} winner for large
vertices; globally it is selected only in its terminal cell. At the
formal vertex zero the two distances tie, but $v=C/[4(H-G)]$ cannot be
zero at finite positive canonical coefficients. In fact, at fixed $C$,
\[
 R(p)-R(-p)=2Cp=2C\sin(\pi/n)>0
\]
at every finite $H-G$. Thus even the unnormalized spectral gap does not
vanish as $H-G\to\infty$; only its value relative to the diverging
quadratic scale tends to zero.

At fixed $n$ and $C>0$, $H-G\to+\infty$ implies $v\downarrow0$.
Since $\beta_{J-1}>0$, every sufficiently large finite $H-G$ lies in
the terminal strict cell, for every residue class. No finite switch is
added at $\chi=\infty$, and no $n\to\infty$ assertion or P06 input
is made. The only reachable negative cosine is the terminal one for
$n\equiv3\pmod4$.

\section{Translation to the original DM parameters}
For $x=x^*>0$ and $s=1+x$,
\[
 H-G=x(h-g/s^2),\qquad C=c/s>0.
\]
Because $x>0$, $H>G\iff h>g/s^2$. In that regime
\[
 v={c/s\over4x(h-g/s^2)}={cs\over4x(hs^2-g)},\qquad
 \chi={4x(hs^2-g)\over cs}.
\]
All factors in the vertex denominator are positive, and no zero
denominator is hidden. For a positive midpoint, solving $v=\beta_j$
gives exactly
\begin{equation}
 h={g\over(1+x)^2}+{c\over4(1+x)x\beta_j}. \label{a05:switch}
\end{equation}
Both terms are positive, so $h>0$ and $H>G$. A nonpositive midpoint
cannot satisfy this concave-regime equation. All activities and cosines
are dimensionless under the contract. The rate coefficients have the
same time scaling; $C,G,H,d$ are rates, $\chi,v,\beta_j,\alpha_j$
are dimensionless, and both terms in \eqref{a05:switch} have the same
units as $h$. No time normalization or hidden calibration is introduced.

\section{Canonical realization and original-equation substitution}
The original synchronous component of the DM vector field is
\[
 f(x)=ax-kx^2+{cx\over1+x}-{gx\over1+x}-hx^2
 =xB(x),\qquad B(x)=a-(k+h)x+{c-g\over1+x}.
\]
For any prescribed $x,c,g,h>0$, the parameter
$a=(k+h)x-(c-g)/(1+x)$ is positive exactly when
$k>(c-g)/[(1+x)x]-h$. In conjunction with $k>0$, this is precisely
\[
 k>\max\{0,(c-g)/[(1+x)x]-h\}.
\]
The right side is finite, so such a $k$ always exists. This equivalence
also covers $c\le g$ and negative lower bounds; no upper bound on $k$
is imposed.

Prescribe any $v>0$ and any $x,c,g>0$, choose $h$ by
\[
 h=g/(1+x)^2+c/[4(1+x)xv],
\]
then choose $k$ as above and define $a$. Every physical coefficient is
strictly positive. Substitution into the \emph{five original terms}
gives
\[
 ax-kx^2+cx/(1+x)-gx/(1+x)-hx^2
 =x[(k+h)x-(c-g)/(1+x)-(k+h)x+(c-g)/(1+x)]=0.
\]
Every component vanishes because the state is synchronous. Moreover
$H-G=c/[4(1+x)v]$ and $C/[4(H-G)]=v$. The construction therefore
realizes the vertex, not merely an unattached spectral polynomial.

To select zero strictly one may take $v=(1+\beta_0)/2$; to select
$1\le j<J$ take $v=(\beta_{j-1}+\beta_j)/2$; to select $J$ take
$v=\beta_{J-1}/2$; to realize switch $j$ take $v=\beta_j$.
All these choices are positive and strict where required. This is a
single universal constructive family, valid for all sizes and parity
classes, not several incompatible existence arguments.

The simpler family in P05 is
\begin{equation}
 x=1,\quad c=g=4,\quad h=1+\chi/2,\quad k>0,\quad
 a=k+h,\quad\chi>-2. \label{a05:family}
\end{equation}
It gives $C=2,G=1,H=1+\chi/2$ and exactly the prescribed $\chi$.
Every positive $\chi$, hence every retained strict cell and switch, is
realized in this family. At $k=1$, $\chi=-1,0,1$ gives respectively
$T>0$, $T=0$, and $T<0$ with $v=1$; all select zero. The domain
$\chi>-2$ ensures $h>0$; the draft never claims this particular
family realizes every negative real value of $\chi$.

\section{Openness, branch coordinates and switch codimension}
Let the equilibrium graph be the subset of positive
$(x,a,k,c,g,h)$ with $B(x)=0$. Since $B_a=1$, it is a smooth
five-dimensional manifold parameterized globally by
\[
 \mathcal U_{\rm eq}=\{(x,k,c,g,h)\in(0,\infty)^5:
 (k+h)x-(c-g)/(1+x)>0\}.
\]
This is open in $\RR^5$. Its projection to physical parameter space is
\[
 \Phi(x,k,c,g,h)=((k+h)x-(c-g)/(1+x),k,c,g,h).
\]
The bottom four rows of its Jacobian are the corresponding identity rows,
so expansion in the first column gives
\[
 \det D\Phi=k+h+{c-g\over(1+x)^2}=-B_x.
\]
At \emph{every} point with $H>G$ this is
\[
 k+(h-g/(1+x)^2)+c/(1+x)^2>0. \label{a05:regularity}
\]
Consequently the scalar root is simple. The original scalar field also
has nonzero derivative $f'(x)=xB_x<0$ at this positive root. There is
no repeated-root/fold point on a retained switch or a concave-regime
cell. This statement requires no full-system Jacobian invertibility.

The inverse function theorem makes $\Phi$ a local diffeomorphism at
these points; equivalently the implicit function theorem solves for a
specified smooth positive $x=x^*(a,k,c,g,h)$. Since $\chi$ is smooth
on the positive equilibrium chart, each strict inequality in
\eqref{a05:cells} gives an open subset of that chart. Each strict cell,
including zero's, has a concave-regime example. Choose a small open
neighborhood within that cell on which $\Phi$ is a diffeomorphism.
Its image is a nonempty open subset of the \emph{full five-dimensional}
positive physical parameter space. Positivity of $x$ and all five
parameters and all strict selection inequalities persists after shrinking
this neighborhood. This proves the claimed physical openness, not
merely openness in a constrained family such as $c=g=4,x=1$.

The quantifier is branch-specific and existential for each cell, not
``every branch at every parameter tuple has the same selected mode''.
For a concrete check of that distinction, $(a,k,c,g,h)=(17/2,1,1,11,1)$
has positive roots $x=1/4$ and $x=3$, by direct substitution. At the
former $\chi=-151/20<0$; at the latter $\chi=15$. For $n=8$ they
select zero and two, respectively. Thus the branch qualification is
necessary and correctly retained. Also $(6,1,1,9,1)$ at $x=1$ gives
$B_x=0$, $\chi=-10$, and a uniform selected representative; there
are canonical folds in the uniform region. P05 explicitly does not
infer its local-diffeomorphism argument at arbitrary such points.
These are scope checks, not proofs of later bifurcation claims.

For $\beta_j>0$ the audit's coefficient-scale switch function is
\[
 \widetilde\Psi_j=x(h-g/(1+x)^2)-{c\over4(1+x)\beta_j}.
\]
Holding branch coordinates $x,k,c,g$ fixed gives
$\partial_h\widetilde\Psi_j=x>0$. P05 uses instead
\[
 \Psi_j=\widetilde\Psi_j/x
 =h-{g\over(1+x)^2}-{c\over4(1+x)x\beta_j},
 \qquad\partial_h\Psi_j=1.
\]
They have the same zero set; division by $x>0$ causes no defect. The
regular-level-set theorem makes it a nonempty four-dimensional smooth
submanifold of the five-dimensional equilibrium chart. Every such point
has $H>G$ by \eqref{a05:switch}, so $\Phi$ is locally regular there.
It transfers the switch to a local codimension-one hypersurface in
$(0,\infty)^5$ physical parameter space on the chosen branch.

For clarity, no fixed-physical-$a$ derivative in $h$ is needed. Holding
$x,k,c,g$ fixed and varying $h$ requires $a'=x$ and gives a legitimate
physical tangent vector $(x,0,0,0,1)$. Along it,
\[
 {\partial\chi\over\partial h}={4(1+x)x\over c}>0,
 \qquad {\partial(R(u_j)-R(u_{j+1}))\over\partial h}
 =-4x(u_j-u_{j+1})\beta_j<0.
\]
Equivalently, for a local inverse chart the physical switch function is
$\Psi_j\circ\Phi^{-1}$, with nonzero differential because $D\Phi^{-1}$
is invertible. The two neighbors exchange order across the surface,
and the finite strict gap to every other representative persists locally.
Thus this is a genuine local change of dominance, not just an equation
with a nonzero derivative on an unrelated space. The draft makes no
unqualified global manifold claim after forgetting branch labels.

\section{Stable paths and what they prove}
For any retained switch choose $\epsilon>0$ smaller than its distance
to zero and to every other threshold, and put
$b_\pm=\alpha_j\pm\epsilon$. When there is no neighboring threshold
on one side, only the remaining positive distances constrain $\epsilon$.
Use \eqref{a05:family} for $\chi\in[b_-,b_+]$, holding
$k>b_+/2$ fixed. All five coefficients are positive, $x=1$ stays an
exact equilibrium, and its scalar derivative is $-(k+h)\ne0$.
The path meets exactly the intended switch and its two neighboring
strict cells, transversely in the selection-surface sense.

Derive the diagonal independently from P04's verified formula
\[
 d=a-2kx-cx/(1+x)^2-g/(1+x)-hx.
\]
With $a=k+h,x=1,c=g=4$, this is
$d=k+h-2k-1-2-h=-k-3$. Equivalently the positive-equilibrium
simplification is $d=-kx-c(1+2x)/(1+x)^2$, used only after imposing
the equilibrium identity. Thus for \emph{every} label and every point
in the entire interval,
\[
 r_m=-k-3+2\cos\theta_m-\frac\chi2\cos(2\theta_m)
 \le-k-1+\frac{b_+}{2}<-1.
\]
The bound uses $\cos\theta\le1$ and $-\cos(2\theta)\le1$ with
$0<\chi\le b_+$. It is uniform in all modes and the full path.
P04 therefore gives strict linear stability throughout. Increasing $k$
while increasing $a$ equally leaves $x,C,G,H,\chi$ fixed in this family
and subtracts the increase of $k$ from every real part. This also
constructs a stable realization of every strict cell. No P11 simplicity,
Hopf theorem, nonzero frequency or instability transversality is used.

What is proved is that dominant-mode switching need not coincide with
loss of stability. It does not identify the first destabilizing mode on
another path, a Hopf event, or any nonlinear bifurcation. The word
``transversely'' in the selection-surface explanation above refers only
to a nonzero switch-function derivative, not to P12's imaginary-axis
crossing condition.

The verified P04 eigenvalue expression is
$\lambda_m=d+Ce^{-i\theta_m}+Ge^{2i\theta_m}-He^{-2i\theta_m}$.
Taking imaginary parts gives independently
\[
 \omega_m=-C\sin\theta_m+(G+H)\sin(2\theta_m).
\]
Hence $r_j=r_{j+1}$ alone is not $\lambda_j=\lambda_{j+1}$;
equality of signed imaginary parts is additionally required. Under
conjugacy these signs reverse. Counts of selected labels are not
counts of distinct eigenvalues or certifications of a simple Hopf pair.
P05 explicitly preserves these distinctions.

\section{Independent small-dimension reconstruction}
Within each row put $\gamma_t=\cos(t\pi/n)$. These are exact symbols,
not rounded constants. The complete representative grids independently
obtained by reflecting angles about $\pi/2$ are
\begin{center}\small
\begin{tabular}{@{}c c l@{}}
\toprule $n$&$N$&$(u_0,\ldots,u_N)$\\\midrule
5&2&$(1,\gamma_2,-\gamma_1)$\\
6&3&$(1,1/2,-1/2,-1)$\\
7&3&$(1,\gamma_2,-\gamma_3,-\gamma_1)$\\
8&4&$(1,\sqrt2/2,0,-\sqrt2/2,-1)$\\
9&4&$(1,\gamma_2,\gamma_4,-1/2,-\gamma_1)$\\
10&5&$(1,\gamma_2,\gamma_4,-\gamma_4,-\gamma_2,-1)$\\
11&5&$(1,\gamma_2,\gamma_4,-\gamma_5,-\gamma_3,-\gamma_1)$\\
12&6&$(1,\sqrt3/2,1/2,0,-1/2,-\sqrt3/2,-1)$\\\bottomrule
\end{tabular}\end{center}
The cells, finite switch boundaries and terminal signs are as follows.
An ordered list in the switch column gives $\alpha_0,\ldots,\alpha_{J-1}$.
\begin{center}\small
\begin{tabular}{@{}c c c c l@{}}
\toprule
$n$&$J$&reachable&$u_J$&finite switch boundaries\\\midrule
5&1&$0,1$&$\sin(\pi/10)>0$&$\gamma_1^{-2}$\\
6&1&$0,1$&$1/2>0$&$4/3$\\
7&2&$0,1,2$&$-\sin(\pi/14)<0$&$\gamma_1^{-2},(\gamma_1\gamma_3)^{-1}$\\
8&2&$0,1,2$&$0$&$4/(2+\sqrt2),2\sqrt2$\\
9&2&$0,1,2$&$\sin(\pi/18)>0$&$\gamma_1^{-2},2/\gamma_1$\\
10&2&$0,1,2$&$\sin(\pi/10)>0$&$\gamma_1^{-2},(\gamma_1\gamma_3)^{-1}$\\
11&3&$0,1,2,3$&$-\sin(\pi/22)<0$&$\gamma_1^{-2},(\gamma_1\gamma_3)^{-1},(\gamma_1\gamma_5)^{-1}$\\
12&3&$0,1,2,3$&$0$&$4/(2+\sqrt3),4/(\sqrt3+1),4$\\\bottomrule
\end{tabular}\end{center}
The positive midpoints, obtained before taking reciprocals, are
\[
\begin{array}{c|l}
5&\gamma_1^2\\
6&3/4\\
7&\gamma_1^2,\ \gamma_1\gamma_3\\
8&(2+\sqrt2)/4,\ \sqrt2/4\\
9&\gamma_1^2,\ \gamma_1/2\\
10&\gamma_1^2,\ \gamma_1\gamma_3\\
11&\gamma_1^2,\ \gamma_1\gamma_3,\ \gamma_1\gamma_5\\
12&(2+\sqrt3)/4,\ (\sqrt3+1)/4,\ 1/4.
\end{array}
\]
For $n=6$ the additional midpoint $\beta_1=0$; for $n=10$ it is
$\beta_2=0$. These are the only zero midpoints in this table and give
no finite switch. In every other row $\beta_J<0$. Monotonicity gives
all remaining signs. Every entry agrees with P05 \lin{487--494};
the four terminal examples in \lin{462--465} also agree exactly.
This small-$n$ verification does not begin P07's separate specialization.

\section{All explicit canonical examples}
Every example in P05 uses \eqref{a05:family} with $k=1$, so its tuple
is $(2+\chi/2,1,4,4,1+\chi/2)$, $x^*=1$, and
\[
 (C,G,H)=(2,1,1+\chi/2),\qquad d=-4,
 \qquad a-k-h+(c-g)/2=0.
\]
All coefficients are positive for the stated $\chi$. Substitution into
the original DM field has already been displayed; here the original
five terms are exactly $a,-1,2,-2,-h$ and sum to zero. For each row
below, substituting every grid value in
\[
 R(u)=-4+\chi/2+2u-\chi u^2
\]
or applying the proved all-grid distance comparison gives precisely
the selected representatives shown.
\begin{center}\small
\begin{tabular}{@{}c l c l@{}}
\toprule
$n$&$\chi$&$S$&maximum real part\\\midrule
any&$-1$&$\{0\}$&$-3/2$\\
any&$0$&$\{0\}$&$-2$\\
any&$1$&$\{0\}$&$-5/2$\\
8&$2$&$\{1\}$&$-4+\sqrt2<0$\\
6&$4/3$&$\{0,1\}$&$-8/3$\\
8&$2\sqrt2$&$\{1,2\}$&$-4+\sqrt2<0$\\
8&$4\sqrt2$&$\{2\}$&$-4+2\sqrt2<0$\\
9&$4/\cos(\pi/9)$&$\{2\}$&$-2+2\sin(\pi/18)<0$\\
10&$2/[\cos(\pi/10)\cos(3\pi/10)]$&$\{2\}$&$7(\sqrt5-5)/10<0$\\
11&$2/[\cos(\pi/11)\cos(5\pi/11)]$&$\{3\}$&$-4-2t+1/t>0$, $t=\sin(\pi/22)$\\\bottomrule
\end{tabular}\end{center}
The $\chi=0$ row is also the tuple $(2,1,4,4,1)$ in \lin{443};
the other three explicitly printed tuples give exactly the stated
$\chi$ values. For $n=8,\chi=2$, the vertex $1/2$ lies between
$\sqrt2/4$ and $(2+\sqrt2)/4$. For the $n=6$ and second $n=8$
examples it equals $\beta_0$ and $\beta_1$, respectively. Their
Fourier labels are $\{0,1,5\}$ and $\{1,2,6,7\}$.

In the last four rows $\chi=2/\beta_{J-1}$, so
$v=\beta_{J-1}/2$ is strictly in the terminal cell for each residue
class. This establishes the maximum against all other labels without
rounding trigonometric constants. The $n=11$ last row is \emph{unstable}
at $k=1$: $0<t<2/11$ from $\sin y<y$ and $\pi<4$, so
$4t+2t^2<96/121<1$, proving the displayed positivity.
P05 asserts this example only as a selection realization, not as a
stable one. Its separate stable-realization assertion explicitly enlarges
$k$. Thus the positive maximum here is not an error or an exception
to the stable-path construction.

\section{Independent computational corroboration and its limits}
The new standard-library checker
\path{A05_support/independent_check.py} does not import P05's validation
code. It computes $\pi$ by Machin's arctangent formula and cosines by
Taylor series at 120 decimal digits, then directly maximizes the
\emph{original P04 trigonometric real part over all $n$ labels}.
It compares that result with the independently stated cell predictions
and also checks the quadratic expression. The tie tolerance is
$10^{-85}\max(1,|r|,|\widetilde r|)$; this is finite-precision
corroboration, not interval arithmetic or a proof of transcendental
equalities. Exact equalities and inequalities are proved above.

The number of finite switches is determined before trusting any draft
statistic:
\[
 \sum_{n=5}^{40}\left\lfloor{n+1\over4}\right\rfloor
 =2\cdot1+4\sum_{r=2}^{9}r+2\cdot10=198.
\]
Indeed $n=5,6$ contribute two ones, $n=7$ through $38$ give four
copies each of $2,\ldots,9$, and $n=39,40$ give two tens.
There are $198+36=234$ strict cells. These are independent analytic
counts, though the switch count agrees with the draft's report.

All \textbf{2,939 independently constructed cases pass}. Coverage is:
\begin{center}
\begin{tabular}{@{}l r@{}}
\toprule Case class&count\\\midrule
Three interior vertices in every strict cell&702\\
Same vertices realized with a second choice of $x,c,g,k$&702\\
Every positive-midpoint tie&198\\
Immediately below / immediately above each switch&198 / 198\\
Stable paths: left, tie, right for every switch&594\\
$T>0$, $T=0$, $v=1$, $\chi=10^{30}$&36 each\\
All negative-midpoint equalities, below the maximum&189\\
Explicit regime / printed examples / terminal examples&3 / 4 / 4\\
Uniform fold / two positive branch cautions&1 / 2\\\bottomrule
\end{tabular}
\end{center}
The adjacent probes are at threshold $\pm10^{-12}$ times a quarter
of the smaller local threshold gap. No nearby second threshold is crossed.
All $5\le n\le40$ and all four residue classes are covered. Every case
uses positive canonical coefficients and verifies the original five-term
DM equilibrium residual. The second construction uses
$x=2/3,c=7/5,g=11/7,k=3$; here $c-g<0$, so $a>0$ is automatic.
The $594$ path sample points satisfy the uniform analytic stability bound;
the proof of the entire intervening intervals is the inequality above.

On Python 3.9.6, the largest scaled original-DM residual was below
$3.18\times10^{-120}$, and the largest scaled direct-versus-quadratic
discrepancy was below $2.45\times10^{-89}$. Full case inputs, selected
labels, maxima and residuals are retained in
\path{independent_results.json}; counts and environment are in
\path{independent_summary.json}. These numerical figures certify only
this newly run checker, not P05's historical run.

\paragraph{U01: historical validation provenance, nonblocking.}
The task does not include \path{P05_validation_support.zip}, which the
manifest says was delivered separately. Its exact symbolic-check count,
1,872 main cases, 2,466 canonical cases, 5,616 dispersion comparisons,
reported error maxima, historical build records and original checker are
not independently reproduced. P05 \lin{558--567} is therefore
\textbf{UNRESOLVED} as a historical provenance assertion. The independent
count $198$ and new checks are not evidence that the absent run itself
occurred as reported. The analytic proof needs none of it. No mathematical
verdict or downstream mathematical gate is blocked by U01; closure should
retain this qualification unless the original companion is supplied and
checked later.

\section{Adversarial counterexample search}
The strongest claims were attacked both analytically and through the
independent all-label comparisons. Failure to find a numerical example
is not used as an impossibility proof. The decisive conclusions are:
\begin{longtable}{@{}L{0.34\textwidth}L{0.61\textwidth}@{}}
\toprule Attempted failure&Outcome and analytic reason\\\midrule\endhead
A finite selected $j>J$&Impossible: its cell has upper endpoint
$\beta_{j-1}\le0<v$; convex cases select only zero.\\
Loss of the negative terminal mode for $n\equiv3\pmod4$&Impossible:
its exact positive upper midpoint leaves a nonempty interval of positive
vertices inside its strict cell.\\
Nonadjacent or triple maximal ties&Impossible: a point strictly between
equidistant endpoints is strictly closer to the vertex.\\
A realizable nonpositive-midpoint maximal switch&Impossible: positive
vertices exclude it; convex/linear regimes have unique representative zero.
All 189 negative-midpoint equalities sampled are lower-value ties.\\
Nonpositive $a$ in the construction&Impossible under the stated strict
$k$ bound, which is algebraically equivalent to $a>0,k>0$.\\
A constructed cell or switch at a scalar fold&Impossible in $H>G$ by
$-B_x=k+(h-g/s^2)+c/s^2>0$. A fold does exist in the uniform convex
region, but P05 excludes the unwarranted universal projection assertion.\\
Vanishing switch gradient after projection&Impossible locally:
$\partial_h\Psi_j=1$ on the chart and $D\Phi$ is invertible.\\
Another mode becoming nonnegative on the stable path&Impossible:
every real part is less than $-1$ on the full interval.\\
A selected Nyquist singleton&Impossible even without $J$:
$R(1)-R(-1)=2C>0$.\\
``Mode zero is always selected''&Counterexample found and exact:
$n=8,(3,1,4,4,2),x=1$ selects $\{1,7\}$. This identifies E01's
wording defect, not a defect in a displayed theorem.\\
All branches at a tuple share selection&Counterexample given above:
$(17/2,1,1,11,1)$ at $n=8$ has different selected representatives on
its two positive roots. P05 does not claim otherwise.\\
All printed realizations are stable at $k=1$&False for the $n=11$
terminal example, but no such universal claim is made. The stable family
chooses $k$ larger explicitly.\\\bottomrule
\end{longtable}

\section{Line-by-line issue ledger for the exact P05 source}
The ranges below partition all 591 lines, including headings, comments
and intervening blank lines. Each substantive range is rated using the
requested vocabulary. PASS means its mathematical or scope assertion
passes; MINOR ISSUE is a required textual correction; UNRESOLVED below
concerns historical provenance only. No MAJOR ISSUE or FALSE theorem
range is present. The literal selected-set reading of E01 is false, but
the correct intended candidate-set meaning is fixed by one local sentence.
\begin{longtable}{@{}L{0.11\textwidth}L{0.18\textwidth}L{0.65\textwidth}@{}}
\toprule Lines&Rating&Independent finding\\\midrule\endhead
1--14&PASS&Current version and DRAFT status; dependency and strengthened scope
metadata agree with the final manifest.\\
15--29&PASS&Positive synchronous hypotheses and verified P04 input, including
the sign $G-H$, are correct.\\
30--31&MINOR ISSUE&E01: candidate set versus selected set is ambiguous.
Require the explicit replacement in section 1. Coprimality sentence passes.\\
32--34&PASS&Only verified earlier modules are imported.\\
35--48&PASS&Global ratio $\chi=-4T/C$ is well defined since $C>0$;
novelty and indexing scope are accurate.\\
49--62&PASS&Quadratic reduction and constant cancellation; no independent
physical freedom of $d$ is inferred.\\
63--83&PASS&Strictly descending grid, both endpoints and disjoint label
classes match the independent derivation.\\
84--94&PASS&Explicit endpoint gap and positivity for every $u<1$ prove
uniqueness, including odd $n$ and $T=0$.\\
95--107&PASS&Concave derivative, positive vertex, $v=1$ uniqueness and
distinction from the first discrete switch.\\
108--120&PASS&Square completion and cosine-space distance are exact.\\
121--129&PASS&Midpoint product; valid range $0\le j<N$; difference formula
for consecutive midpoints applies where $j<N-1$.\\
130--145&PASS&Theorem A's closed endpoint/interior cells and unique open
cells, ties and positive-midpoint restriction are correct.\\
146--163&PASS&Distance-difference sign and neighbor sufficiency against
all other representatives are correct.\\
164--171&PASS&Nonadjacent/triple maximal tie exclusion; independent explicit
distance proof supplied in this audit.\\
172--184&PASS&Adjacent gap $C(u_j-u_{j+1})(1-\chi\beta_j)$; nonmaximal
convex equalities are correctly distinguished.\\
185--193&PASS&Three/four selected-label counts; no distinct-eigenvalue or
Nyquist-selection inference.\\
194--210&PASS&Theorem B's exact list, strict attainability and finite
positive switches are all justified.\\
211--231&PASS&Cell-intersection criterion, integer floor formula,
$J<N$, and no isolated additional reachable representative.\\
232--245&PASS&All four rows of the parity table, including formal zero
ties, are correct.\\
246--262&PASS&Odd straddling cosines, sine ordering, zero-cosine pair and
zero-midpoint exception.\\
263--280&PASS&Fixed-$n$ terminal limit and persistent $2C\sin(\pi/n)$
gap; no finite infinite-$\chi$ switch or large-$n$ claim.\\
281--300&PASS&Reciprocal directions, ranges of $\alpha_j$, combined
global $\chi$ rule and strict uniform status at $\chi=1$.\\
301--318&PASS&Canonical inequalities, positive denominators and exact
switch equation.\\
319--328&PASS&Original equilibrium constraint and exclusion of
nonpositive-midpoint dominant switches.\\
329--341&PASS&Theorem C quantifiers are existential for open cells and
local on specified branches for hypersurfaces; all switch points are regular.\\
342--362&PASS&Universal construction, exact strict $k$ bound and original
five-term residual; spectral coefficients are not treated as free.\\
363--376&PASS&Simple family has $\chi>-2$ and covers all positive
$\chi$; $-1,0,1$ examples are valid.\\
377--395&PASS&Open five-dimensional equilibrium chart and correctly signed
projection determinant.\\
396--403&PASS&Strict positivity of the determinant throughout $H>G$;
IFT hypotheses and scalar versus full-system regularity are distinguished.\\
404--413&PASS&Open strict cells project locally to open physical sets;
concave uniform example included; arbitrary uniform folds are not claimed regular.\\
414--425&PASS&Scaled switch function has derivative 1, nonempty regular
level set and valid local transfer to physical parameters.\\
426--439&PASS&Both crossing derivatives correct; $a$ varies to hold $x$
fixed; other representatives remain strictly below; no global branch quotient claim.\\
440--453&PASS&Four printed tuples have positive coefficients, $x=1$, exact
residuals and stated label maxima.\\
454--471&PASS&All four terminal examples use $v=\beta_{J-1}/2$;
the unstable $n=11,k=1$ example makes no stability claim.\\
472--479&PASS&Exact row-local notation, terminal midpoint signs and P07
scope restriction.\\
480--498&PASS&Every grid and midpoint entry for $n=5,\ldots,12$ agrees
with independent reconstruction.\\
499--507&PASS&Dominance and oriented first loss of stability are separated;
no first-instability result is inferred.\\
508--520&PASS&Stable path remains canonical, $d=-k-3$ and uniform
all-label bound $r_m<-1$ hold throughout the interval.\\
521--527&PASS&Adjacent change of dominance and stable realization of every
cell; later first-crossing tasks are explicitly deferred.\\
528--541&PASS&Imaginary-part sign correct; real-part ties do not establish
complex equality, simple pairs or Hopf conditions.\\
542--556&PASS&Internal analytic checks describe true identities; the
audit independently proves them rather than trusting their description.\\
557--563&UNRESOLVED&U01: contents and coverage of the separately named
original validation companion cannot be authenticated from this delivery.\\
564--565&PASS&Finite-precision corroboration is explicitly distinguished
from exact trigonometric input definitions and proof.\\
566--567&UNRESOLVED&U01: historical support run totals, environment and
error records are absent; the new checker is not that run.\\
568--569&PASS&Compilation and internal checking do not confer verification.\\
570--591&PASS&Auditor checklist is consistent with the verified findings and
scope; it is source material, not a new instruction to alter any proof.\\
\bottomrule
\end{longtable}

\section{Dependencies, strengthened target and downstream consequences}
The audit uses only the supplied verified P01--P04 input, elementary
trigonometry, finite-set comparisons, differentiation, and the local
implicit/inverse function and regular-level-set theorems with their
nonzero derivatives explicitly checked. No result from P06, P07, P08,
P11, P12 or P13 is a premise. Later module names in P05 are scope
deferrals, not circular dependencies. No new literature search is conducted.

The mathematical content reaches the frozen strengthened P05 target:
an exact finite-$n$ classification of all reachable cells and ties, the
modulo-four terminal phenomenon and excluded Nyquist mode, exact original
parameter realization, open physical cells and realizable local switch
geometry, plus stable crossing paths. The bare quadratic maximizer is
only its elementary tool. E01 does not reduce this mathematical content;
it requires that the introductory definition say what the later theorems
actually prove. This audit does not certify new-model priority, an exact
DM--MB contrast, or the later first-Hopf package. The frozen literature
priority limitation remains separate.

\textbf{P08 clearance:} after E01 is applied in a separately authorized
closure and that closure records the exact corrected version, the P05
mathematical prerequisite for P08 is cleared. A closure that merely
changes status without correcting E01 is insufficient. Retain U01 as
a nonblocking provenance note unless independently resolved. P08 must
still respect the frozen same-coefficient tuple, original time variable,
and separately specified positive branches on their common existence
domain; P05 does not prove a comparator contrast. No P08 proof is begun.

\textbf{P07 consequence:} after the same closure, the general selection
theorem may be specialized. The small-$n$ checks here neither begin nor
promote the separate five-node theorem. \textbf{P11 consequence:} exact
dominance cells and real-part ties may be used after closure, but branch,
physical path, frequency, simple critical pair, all-other-mode exclusion
and first-crossing conditions remain to be proved there and in its
successors. The stable paths explicitly show why a selection switch
alone cannot supply them. No P11 proof is begun.

\paragraph{Final disposition.}
Keep the supplied P05 file, project manifest, claim registry and compendium
unchanged. Deliver this standalone audit, its independently written
corroboration and hash records as an audit-only package. The current
source still says DRAFT. Required correction: E01 only. The analytic
verdicts are independent of both the historical and new numerical checks.

\appendix
\section{Exact dependency and context hashes}
All hashes in this appendix are full SHA-256 digests of supplied bytes. Body hashes use the first line-anchored section command through EOF. The accompanying JSON preserves byte and line counts as well.
\medskip\noindent\textbf{\path{project_manifest.md}}
86494 bytes; 451 lines. Full file:
\sha{a4658ef7b17b64fbb43825bfa1cd035d}{f07efa1fb72141f9c9fea02dd1c28830}
\medskip\noindent\textbf{\path{claim_registry.md}}
92791 bytes; 465 lines. Full file:
\sha{12c106cf55d789d840329a2f5f4e80e1}{a54163707d4d4286fd387bbefeb149ee}
\medskip\noindent\textbf{\path{frozen_theorem_package.md}}
9729 bytes; 69 lines. Full file:
\sha{724adf3d04a7f9661a2e53fbae359192}{e8a5a5ee56737bdca2379209b556ea63}
\medskip\noindent\textbf{\path{final_novelty_verdict.md}}
6511 bytes; 37 lines. Full file:
\sha{5fa8ddd5ada9628b104dba68dd535356}{442c5e62ec45b72f3c2c354dec6037ea}
\medskip\noindent\textbf{\path{novelty_decisions.md}}
19241 bytes; 124 lines. Full file:
\sha{8fd079fb108b93c4504bc2b527de11c1}{f7e867a967c7cef0d90e1ff7a0a9175e}
\medskip\noindent\textbf{\path{02_theorem_dependency_map.tex}}
9372 bytes; 145 lines. Full file:
\sha{a1a4435f36fc1d1b6f0b2c5226a91328}{3c76658d55b96493e6c50a8693831cac}
\medskip\noindent\textbf{\path{00_canonical_model.tex}}
3712 bytes; 74 lines. Full file:
\sha{c209981cb2924e8a782f24aa2904e8c8}{5a105ff99319dd6aada12f48a351c2a4}
\medskip\noindent\textbf{\path{01_notation_and_assumptions.tex}}
5152 bytes; 103 lines. Full file:
\sha{97d993c6d3986524dbcc29659ee73d25}{7be1b03a35c44401e8ef4534303eeb91}
\medskip\noindent\textbf{\path{proofs_compendium.tex}}
5807 bytes; 124 lines. Full file:
\sha{d1b9c83900c027edffa9c1fde6088546}{c43cd37750e602636668b0abd53eb09a}
\medskip\noindent\textbf{\path{shared_macros.tex}}
1281 bytes; 35 lines. Full file:
\sha{f4a18a440e1214b619bf18f4a81fee50}{7239460ca796d8ca3916791069d6a7ad}
\medskip\noindent\textbf{\path{proofs/P01_wellposedness_and_absorbing_box.tex}}
22570 bytes; 520 lines. Full file:
\sha{0d6950821aa83f3d09e8bfe84dab2b0b}{5020aa6e7881cf6928e984bff1890f1e}
Mathematical body:
\sha{a640b2bf9a160feed273dc667c6b44a6}{923a1d68d6dd81acd35640cc0e4d5808}
\medskip\noindent\textbf{\path{audits/A01_wellposedness_and_absorbing_box.tex}}
48622 bytes; 974 lines. Full file:
\sha{8d5c1b2ec782dbe49b07d2bbe529ee47}{4d4c142695c53b523ccb41d63bb33098}
\medskip\noindent\textbf{\path{proofs/P02_synchronous_subspace_and_equilibria.tex}}
26905 bytes; 610 lines. Full file:
\sha{8e057c2cd009659eee45906844e25f4f}{38c156e5eb2e4130e824aa5e8299c51a}
Mathematical body:
\sha{8a061ba81c7ffd2d05b188b9fca4ec1b}{94bc132147f7f664c8523f0a5f32560c}
\medskip\noindent\textbf{\path{audits/A02_synchronous_subspace_and_equilibria.tex}}
37895 bytes; 847 lines. Full file:
\sha{ca3d723d9fcfac14f5129583441e611a}{65e41848ec44f3b5e3a0072f795d84ac}
\medskip\noindent\textbf{\path{proofs/P03_origin_and_linearization.tex}}
19857 bytes; 447 lines. Full file:
\sha{6c3c9bfb2307bdc043eb7f6a874cb222}{e2ac2ce0750b339181f0088dd0bd859e}
Mathematical body:
\sha{07d76c8ce0b5232827e9f2c74245188b}{55b0ce0d1732b3cb2e528d69e8d6b50a}
\medskip\noindent\textbf{\path{audits/A03_origin_and_linearization.tex}}
37700 bytes; 692 lines. Full file:
\sha{e1008f98ba9667b92b8d74f1bc7d21e0}{21414b454ac976ae0db1500b4b7ef8d1}
\medskip\noindent\textbf{\path{proofs/P04_arbitrary_nq_fourier_spectrum.tex}}
26980 bytes; 590 lines. Full file:
\sha{1331c8598acd0481046250d918cac028}{ca6d13d2604612448c3459e864321c58}
Mathematical body:
\sha{3e3d58230c87af3b5f07069d61065362}{dc081d172ad3b236057fefda25fa5653}
\medskip\noindent\textbf{\path{audits/A04_arbitrary_nq_fourier_spectrum.tex}}
47595 bytes; 832 lines. Full file:
\sha{8a57079241269c92a3c3e8dc4dccf859}{2da2d1a636fca0c0990a4f318396d869}
\medskip\noindent\textbf{\path{proofs/P05_q2_mode_selection.tex}}
26848 bytes; 591 lines. Full file:
\sha{b7123736242554bd7fc1f6a56d8788b6}{c5413a25930e5360d7ea5babe44184d4}
Mathematical body:
\sha{25ab3e4ed5d871ab2884210e6d8e7f36}{31acee3f7a1c662d3c33de441048810d}
\end{document}
